\clearpage
\begin{summarybox}
  \textbf{\textcolor{CSummary}{核心一句话}}
  \textbf{轴对称：横坐标取平均，纵坐标相等。}
  对称轴是竖直线 $x=a$；两横坐标的平均数是 $a$，对应函数值相等：
  \[
    \boxed{f(a+t)=f(a-t)}\qquad\Longleftrightarrow\qquad
    \boxed{f(x)=f(2a-x)}.
  \]
  这里的两种写法等价；要能表示\textbf{整幅图象}的对称，前提是定义域左右成对，等式对全部配对成立。
  看见 $f(p+t)=f(q-t)$ 或 $f(p-t)=f(q+t)$，先算两个输入的和 $p+q$，再读出轴 $x=(p+q)/2$；仍须满足上述全域配对条件。

  \medskip
  \textbf{\textcolor{CSummary}{为什么能记住这个公式？}}
  对称轴 $x=a$ 是两个横坐标的\textbf{中间位置}：
  \[
    \frac{x+(2a-x)}2=a.
  \]
  因此同一对位置的横坐标之和必为 $2a$，对应函数值必相等。

  \medskip
  \textbf{\textcolor{CSummary}{做题三步}}
  \begin{enumerate}[leftmargin=*,itemsep=0.35em,topsep=0.3em]
    \item \textbf{先看范围：}左右两侧的位置能不能都代入？
    \item \textbf{再找配对：}以 $a$ 为中点，另一个位置是 $2a-x$。
    \item \textbf{最后比高度：}证明对所有位置都相等，或利用相等直接求值。
  \end{enumerate}

  \medskip
  \textbf{\textcolor{CSummary}{最值得记住的三个联系}}
  \begin{itemize}[leftmargin=*,itemsep=0.35em,topsep=0.3em]
    \item \textbf{二次函数：}配方可直接读出对称轴。
    \item \textbf{偶函数：}把轴 $x=a$ 移到 $0$ 后，得到 $g(t)=g(-t)$。
    \item \textbf{周期：}同一图象有两条竖直对称轴 $x=a$、$x=b$（$a\ne b$），
      就有周期 $2|b-a|$，但不一定是最小正周期。
  \end{itemize}

  \medskip
  \textbf{\textcolor{CSummary}{最后提醒}}
  \textbf{几对相等只是一点线索，所有位置成对相等才是证明。}
  定义域缺一半，或者只验几组数，都不能直接断定整个图象对称。
\end{summarybox}
